\documentclass[12pt,a4paper]{article}
\usepackage[utf8]{inputenc}
\usepackage[vietnamese]{babel}
\usepackage{graphicx}
\usepackage{listings}
\usepackage{xcolor}
\usepackage{geometry}
\usepackage{hyperref}
\usepackage{float}

\geometry{margin=2.5cm}

\lstset{
    basicstyle=\ttfamily\small,
    frame=single,
    breaklines=true,
    keywordstyle=\color{blue},
    commentstyle=\color{gray},
    stringstyle=\color{red},
    language=Java
}

\title{\textbf{BÁO CÁO BÀI TẬP 01} \\ Lập trình hướng đối tượng}
\author{
    Họ tên: Lương Quốc Khánh \\
    MSSV: 20235754 \\
    Nhóm: 15
}
\date{Tháng 3, 2026}

\begin{document}

\maketitle
\tableofcontents
\newpage

% ============================================================
\section{HelloWorld.java}
% ============================================================

Chương trình yêu cầu người dùng nhập tên, sau đó hiển thị lời chào.

\begin{lstlisting}
import java.util.Scanner;

public class HelloWorld {
    public static void main(String[] args) {
        Scanner sc = new Scanner(System.in);
        System.out.print("Nhap ten cua ban: ");
        String ten = sc.nextLine();
        System.out.println("Xin chao, " + ten + "!");
        sc.close();
    }
}
\end{lstlisting}

% ============================================================
\section{MVCTutorial}
% ============================================================

Ứng dụng Calculator đơn giản sử dụng mô hình MVC với Java Swing. Gồm 4 file:

\subsection{CalculatorModel.java}
Thực hiện các phép tính cộng, trừ, nhân, chia. Không biết sự tồn tại của View.

\subsection{CalculatorView.java}
Giao diện người dùng dùng Swing (JFrame, JTextField, JComboBox, JButton). Chỉ hiển thị, không tính toán.

\subsection{CalculatorController.java}
Điều phối giữa Model và View. Lắng nghe sự kiện click nút ``Tinh'', lấy dữ liệu từ View, gọi Model tính toán, rồi gửi kết quả về View.

\subsection{MVCCalculator.java}
File main khởi tạo Model, View, Controller và hiển thị cửa sổ.

% ============================================================
\section{MVC Calculator - Swing}
% ============================================================

Phiên bản Calculator riêng dùng Swing, cấu trúc tương tự MVCTutorial nhưng đặt tên class có hậu tố ``Swing''. Gồm 4 file:

\begin{itemize}
    \item \texttt{CalculatorModelSwing.java} -- Model xử lý phép tính
    \item \texttt{CalculatorViewSwing.java} -- Giao diện Swing
    \item \texttt{CalculatorControllerSwing.java} -- Controller điều phối
    \item \texttt{MVCCalculatorSwing.java} -- Main
\end{itemize}

Các điểm khác biệt: sử dụng tên biến tiếng Việt (\texttt{txtSo1}, \texttt{cboPhepTinh}, \texttt{btnTinh}), thông báo lỗi bằng tiếng Việt.

% ============================================================
\section{MVC Calculator - JavaFX}
% ============================================================

Phiên bản Calculator dùng JavaFX với FXML (Scene Builder). Dự án Maven gồm:

\begin{itemize}
    \item \texttt{pom.xml} -- Cấu hình Maven, dependency JavaFX 21.0.2
    \item \texttt{Main.java} -- Khởi tạo JavaFX Application, load FXML
    \item \texttt{CalculatorModelJFX.java} -- Model xử lý phép tính
    \item \texttt{CalculatorControllerJFX.java} -- Controller dùng annotation \texttt{@FXML}
    \item \texttt{calculator-view.fxml} -- Giao diện được thiết kế bằng FXML
\end{itemize}

Trong JavaFX, View được định nghĩa trong file FXML, Controller sử dụng \texttt{@FXML} annotation để bind các thành phần UI. Model vẫn giữ chức năng xử lý logic tính toán.

% ============================================================
\section{Các thao tác Git}
% ============================================================

Dưới đây là minh chứng các thao tác Git đã thực hiện.

% ---------- Git Add ----------
\subsection{Git Add}
Thêm các file bài tập vào staging area.

\begin{lstlisting}[language=bash]
git add Homework01/20235754-LuongQuocKhanh/
\end{lstlisting}

\begin{figure}[H]
    \centering
    % \includegraphics[width=0.9\textwidth]{screenshots/git-add.png}
    \fbox{\parbox{0.8\textwidth}{\centering\vspace{3cm}\textit{[Screenshot: git add]}\vspace{3cm}}}
    \caption{Git Add}
\end{figure}

% ---------- Git Commit ----------
\subsection{Git Commit}
Commit các file đã thêm.

\begin{lstlisting}[language=bash]
git commit -m "Add homework files for 20235754-LuongQuocKhanh"
\end{lstlisting}

\begin{figure}[H]
    \centering
    % \includegraphics[width=0.9\textwidth]{screenshots/git-commit.png}
    \fbox{\parbox{0.8\textwidth}{\centering\vspace{3cm}\textit{[Screenshot: git commit]}\vspace{3cm}}}
    \caption{Git Commit}
\end{figure}

% ---------- Git Push ----------
\subsection{Git Push}
Đẩy code lên remote repository.

\begin{lstlisting}[language=bash]
git push origin main
\end{lstlisting}

\begin{figure}[H]
    \centering
    % \includegraphics[width=0.9\textwidth]{screenshots/git-push.png}
    \fbox{\parbox{0.8\textwidth}{\centering\vspace{3cm}\textit{[Screenshot: git push]}\vspace{3cm}}}
    \caption{Git Push}
\end{figure}

% ---------- Modify File ----------
\subsection{Modify File}
Sửa file HelloWorld.java rồi commit lại.

\begin{lstlisting}[language=bash]
git add Homework01/20235754-LuongQuocKhanh/HelloWorld.java
git commit -m "Modify HelloWorld.java"
git push origin main
\end{lstlisting}

\begin{figure}[H]
    \centering
    % \includegraphics[width=0.9\textwidth]{screenshots/git-modify.png}
    \fbox{\parbox{0.8\textwidth}{\centering\vspace{3cm}\textit{[Screenshot: git modify]}\vspace{3cm}}}
    \caption{Modify File}
\end{figure}

% ---------- Git Pull ----------
\subsection{Git Pull}
Kéo code mới từ remote.

\begin{lstlisting}[language=bash]
git pull origin main
\end{lstlisting}

\begin{figure}[H]
    \centering
    % \includegraphics[width=0.9\textwidth]{screenshots/git-pull.png}
    \fbox{\parbox{0.8\textwidth}{\centering\vspace{3cm}\textit{[Screenshot: git pull]}\vspace{3cm}}}
    \caption{Git Pull}
\end{figure}

% ---------- Remove File ----------
\subsection{Remove File}
Xóa file bằng git rm rồi commit.

\begin{lstlisting}[language=bash]
git rm Homework01/20235754-LuongQuocKhanh/temp.txt
git commit -m "Remove temp file"
git push origin main
\end{lstlisting}

\begin{figure}[H]
    \centering
    % \includegraphics[width=0.9\textwidth]{screenshots/git-remove.png}
    \fbox{\parbox{0.8\textwidth}{\centering\vspace{3cm}\textit{[Screenshot: git rm]}\vspace{3cm}}}
    \caption{Remove File}
\end{figure}

% ---------- Create Branch ----------
\subsection{Create Branch}
Tạo nhánh mới.

\begin{lstlisting}[language=bash]
git checkout -b feature-khanh
\end{lstlisting}

\begin{figure}[H]
    \centering
    % \includegraphics[width=0.9\textwidth]{screenshots/git-branch.png}
    \fbox{\parbox{0.8\textwidth}{\centering\vspace{3cm}\textit{[Screenshot: git checkout -b]}\vspace{3cm}}}
    \caption{Create Branch}
\end{figure}

% ---------- Merge Branch ----------
\subsection{Merge Branch}
Merge nhánh feature-khanh vào main.

\begin{lstlisting}[language=bash]
git checkout main
git merge feature-khanh
git push origin main
\end{lstlisting}

\begin{figure}[H]
    \centering
    % \includegraphics[width=0.9\textwidth]{screenshots/git-merge.png}
    \fbox{\parbox{0.8\textwidth}{\centering\vspace{3cm}\textit{[Screenshot: git merge]}\vspace{3cm}}}
    \caption{Merge Branch}
\end{figure}

% ---------- Resolve Conflict ----------
\subsection{Merge - Resolve Conflict}
Tạo conflict giữa 2 branch rồi giải quyết.

\begin{lstlisting}[language=bash]
# Tao branch conflict-test, sua file
git checkout -b conflict-test
# Sua MVCCalculator.java, commit

# Quay lai main, sua cung file do
git checkout main
# Sua MVCCalculator.java, commit

# Merge -> conflict
git merge conflict-test
# Resolve conflict trong file
git add .
git commit -m "Resolve merge conflict"
\end{lstlisting}

\begin{figure}[H]
    \centering
    % \includegraphics[width=0.9\textwidth]{screenshots/git-conflict.png}
    \fbox{\parbox{0.8\textwidth}{\centering\vspace{3cm}\textit{[Screenshot: merge conflict]}\vspace{3cm}}}
    \caption{Merge Conflict}
\end{figure}

\begin{figure}[H]
    \centering
    % \includegraphics[width=0.9\textwidth]{screenshots/git-resolve.png}
    \fbox{\parbox{0.8\textwidth}{\centering\vspace{3cm}\textit{[Screenshot: resolve conflict]}\vspace{3cm}}}
    \caption{Resolve Conflict}
\end{figure}

% ============================================================
\section{Kết luận}
% ============================================================

Bài tập đã hoàn thành các yêu cầu:
\begin{itemize}
    \item Viết chương trình HelloWorld.java
    \item Implement MVCTutorial example
    \item Implement Calculator theo mô hình MVC với 2 phiên bản: Swing và JavaFX (FXML)
    \item Thực hiện đầy đủ các thao tác Git: add, remove, modify, commit, push, pull, create branch, merge branch, resolve conflict
\end{itemize}

\end{document}
